\scanpage{108}
variable dans une multiplicité topologique convenable $M_G^C$, les $\unsure{P_{\mathcal{X}}}$ forment un $\mathbb{U}_{M_G^C}^{C}$-torseur sur $M_G^C$, soit $\widetilde{M}_G^C$, et on définit un \deleted{application} morphisme de Topos
\[
\widetilde{M}_G^C\longrightarrow M_{G'}.
\]
\deleted{$\unsure{M_G}$}
\margin{c'est l'\unsure{image} inverse d'un torseur évident \deleted{$\unsure{P_{G,C}}$ \unsure{sur} $\unsure{M_G}$} $P_{G,C}$ ou $P_C$ sur $\unsure{M_G}$.}

\unsure{Revenant|Revenons} à $\mathcal{X}=(X,(D_{\unsure{\widetilde x|\widetilde a}}))$, \deleted{on} on va définir, pour tout $t\unsure{=|\in}(t_x)_{x\in C}\in T_C$, une \deleted{\unsure{courbe}} $X'$ du type $G'$, notée $X'_t$. On l'obtient par chirurgie holomorphe : partir de $X_1$, en recollant \inserted{$\forall x\in C$} les deux composantes \deleted{$\unsure{\partial_a}$} $T_{x'},T_{x''}$ du bord de $X_1$ correspondant à $x$, suivant l'antiisomorphisme $t_x$. Pour $\mathcal{X},t$ variable, cela donne \unsure{l'homomorphisme|le bon morphisme} de topos ci-dessus.

Soit $M_G\deleted{_{\unsure{G'}}}$ \deleted{\unsure{ou} $\unsure{M_{G,e}}$} la multiplicité holomorphe des courbes $X$ du type $G$, on a donc
\[
M_G^C\longrightarrow M_G\qquad\text{équiv. d'homotopie}
\]
donc
\diagram{\begin{tikzcd}[column sep=small,row sep=small]
1 \arrow[r,dash] & \unsure{\Pi_1|\pi_1}(\mathbb{U}^{C}) \arrow[r] \arrow[d,equal] & \unsure{\Pi_1|\pi_1}(\widetilde{M}_G^C) \arrow[r] \arrow[dd] & \unsure{\Pi_1|\pi_1}(M_G) \arrow[r,dash] \arrow[d,equal] & 1 \\
& \mathbb{Z}^{C} & & \unsure{\Pi_1|\pi_1}(M_G^C) & \\
& & \unsure{\Pi_1|\pi_1}(M_{G'}). & &
\end{tikzcd}}

Pour avoir $S\mathcal{T}_G\longrightarrow S\mathcal{T}_{G'}$, on considère le fibré $SP_{G'}$ sur $M_{G'}$, principal de groupe $\mathbb{U}^{\widehat{A}'}$
\mspage{107}
